\scanpage{89}
\mspage{\handwritten{3\deleted{\unsure{4|7}}8}}
% A single diagonal crosses the typescript, including its footnote.
\deleted{au cours de cette année-là que j'ai appris pour la première fois (mais non \inserted{\handwritten{certes pour}} la dernière\handwritten{!}) \deleted{\unsure{certes|censés} \unsure{!|(*}}\deleted{)} que\deleted{\unsure{la}} c'est considéré comme un manque de sérieux de démontrer bel et bien des choses délicates et visiblement fondamentales, \inserted{\handwritten{que tout le monde utilise et}} que les prédécesseurs se sont toujours contentés d'admettre\inserted{\handwritten{(*)}}. Ou de démontrer un résultat qui englobe comme cas particuliers plusieurs résultats \inserted{profonds} connus (ce qui montre bien évidemment que le résultat soi disant nouveau ne peut être qu'un cas particulier ou une conséquence facil\unsure{e} de\deleted{\unsure{s}}s résultats connus). Ou de prendre la peine, \deleted{d'énoncer un}\inserted{dans \deleted{\unsure{un}}\handwritten{l'énoncé d'un}} résultat ou dans la description d'une situation en termes d'une autre, de formuler avec soin les hypothèses naturelles (signe d'un regrettable bombinage), plutôt que de se \deleted{contenter} borner \deleted{au cas simple} à quelque cas d'espèce au goût de celui qui émet son opinion. Ainsi j'ai vu faire reproche à Contou-Carrère de ne pas s'être borné dans sa thèse à se placer sur le corps des complexes ou tout au moins sur un corps, au lieu d'un schéma général\deleted{\unsure{.}} (tout en lui concédant \inserted{quand même} la circonstance atténuante que \deleted{\unsure{s}}\handwritten{c}'était \deleted{son patron \unsure{qui}} \inserted{\handwritten{sûrement sur les instances de son patron qu'il avait dû s'y résoudre)}}\handwritten{.} Celui qui s'exprimait ainsi était pourtant suffisamment dans le \deleted{\handwritten{\unsure{assister à quelque chose}}} coup pour savoir que même en se bornant au corps des complexes, les nécessités de\handwritten{s} \deleted{\unsure{la}} démonstration\handwritten{s} forçaient la main pour introduire des schémas de base générau\deleted{\unsure{n}}x ...\deleted{\unsure{?!|)}}}
\margin{\handwritten{\tikz[baseline]{\draw (0,0)--(.34,0)--(.40,.28)--(.73,.28); \draw (.31,0)--(.37,.28);}}}

\deleted{Les égarements d'une certaine mode \inserted{\handwritten{aujourd'hui}} vont jusqu'à \deleted{\unsure{mettre} exclure \unsure{tain|tair|tard|Les}}\inserted{honnir} non seulement \deleted{\unsure{d'une}}\inserted{\handwritten{\unsure{d|l}}es} démonstrations\handwritten{\ }en \deleted{bonne et due} forme, mais \inserted{\handwritten{souvent}} même \handwritten{d}es énoncés \deleted{en bonne ou en forme} et \handwritten{d}es définitions en forme\unsure{;|.|:} \deleted{\unsure{x}} au prix \deleted{d'un énoncé en bonne et due forme (sans même parler des hypothèses les plus naturelles), et} où est le papier et la longanimité du lecteur, il ne sera bientôt plus question de s'encombrer d'un luxe aussi coûteux\handwritten{!} \deleted{et} \deleted{on se bornera prochainement à}\inserted{\unsure{En} extrapolant les tendances actuelles}, il ne doit pas être difficile de \deleted{calculer}\inserted{\handwritten{prédire}} le moment exact où il ne sera plus question \inserted{dans une publication} d'expliciter \deleted{énoncés ni} définitions \inserted{\handwritten{ni énoncés}}, qu'on se contentera \deleted{de désigner}\inserted{désormais de nommer} par des mots-code, en laissant à l'infatigable et génial lecteur le soin de remplir les blancs, conformément à ses propres lumières. La tâ\deleted{\unsure{c}}che du referee sera facilitée d'autant, car il lui suffira de \deleted{consulter}\inserted{\handwritten{regarder}} dans l'annuaire \handwritten{``}Wo is Who\handwritten{''} si l'auteur est connu comme crédible (de toutes façons personne ne pourrait contredire les blancs et les \deleted{petits points}\inserted{\handwritten{pointillés}} qui composent le brillant article), ou au contraire un inavouable inconnu\handwritten{\ }qui sera (comme c'est le cas déjà dès aujourd'hui) \deleted{éliminé}\inserted{\handwritten{éjecté}} d'office ...}

\noindent\rule{2em}{.4pt}\par
% slanted
\margin{\handwritten{\unsure{sans interligne|sans interlignes}}}
\deleted{(\textsuperscript{*}) Voir aussi à ce sujet l'épisode ``La note - ou la nouvelle éthique'' \handwritten{(}par\unsure{.}33 \handwritten{)}. \deleted{Qui}\inserted{\handwritten{Cette fameuse ``note''}} avait justement le tort d'expliciter des notions \deleted{ou}\inserted{\handwritten{et}} des \deleted{dém} énoncés \deleted{et des démonstrations} qui avaient été laissées jusque là dans le vague, et qui pourtant ont été \deleted{utilisés (par moi en \unsure{tout autant|tout allant} par \unsure{les|le} \unsure{élèves|élevé|évidé}) sans vergogne}\inserted{implicitement utilisés par moi} pour établir des résultats qui portent mon nom et \deleted{\unsure{que}} que tout le monde \deleted{(y compris \inserted{\handwritten{\unsure{mon}}}\unsure{maître et élève})} utilise sans vergogne\deleted{\unsure{.|,}}\inserted{\handwritten{depuis bientôt vingt-cinq ans.}}}
% archive sheet number: 88
